\begingroup
\csname tagsleft@true\endcsname
\cotefile{156-6}
\scanpage{16}\mspage{15}
\unsure{par|pr} contre la relation $\ll$ n'a pas d'interprétation simple directe, ni en termes des \inserted{\unsure{déployements|déployement}} $\widetilde F\subset\unsure{M|\mathcal{M}}$, ni en termes des multidéployement $\mathrm{Fig}(F)$; ainsi $G\ll F$ n'implique nullement $\mathrm{Fig}_{M}(G)\ll\mathrm{Fig}_{M}(F)$. Aussi l'application ($**$) et la notation $=\mathrm{Fig}_{M}(F)$ sont-\unsure{il|ils} finalement un peu bidon !

En termes de $M$, la structure est donnée par
\[
\begin{array}{|l}\hline
\unsure{\mathcal{F}|\mathfrak{F}|\mathscr{F}}\subset\unsure{\mathcal{P}|\mathscr{P}|\mathfrak{P}|\mathrm{P}}(M)\quad\text{ens de parties de }M\\
\ll\quad\text{relation d'ordre dans }\mathcal{F},\ \text{\uline{impliquée}}\\
\qquad\text{par l'inclusion }(\deleted{\unsure{F}}\text{\unsure{notée|noté} }F\leq G)
\\\hline\end{array}
\tag*{(1.33)}
\]
Ça a l'air moins sympathique du point de vue simplicité des données et d'autant plus qu'il n'y a pas \inserted{à première vue} d'interprétation directe de $\ll$ en termes de \deleted{\unsure{Fig-structure|Fig, structure|notions de structure|notions de théorie}} \inserted{notions ensemblistes} familières. Mais c'est ici que vont servir les (grandes) ombres !

En termes de la \inserted{\unsure{première|premier}} donnée (1.33.\unsure{)|]} \margin{(et avant d'introduire $\ll$)}, le \deleted{\unsure{tout}} premier axiome est \deleted{\illegible{1}} le fait que $\mathcal{F}$ est \deleted{une figure} \inserted{une famille d'}ens. \inserted{\uuline{pré}admissibles} \deleted{\unsure{dans}} $M$ de \margin{attention au ``\uuline{pré}'' !} \deleted{\unsure{Axiome|Axiomes}} support $M$, \deleted{\unsure{d'où}} donc la donnée de \margin{\deleted{\unsure{un ens.|un ensemble|Chaque ens} de \unsure{parties} \unsure{fermées|finies} de $\unsure{F|f|\mathcal{F}|\lbrace}$}} $\mathcal{F}$ \deleted{\unsure{revient à|revient} la donnée d'} \inserted{implique} une relation de préordre dans $M$ (\unsure{ce sera|ce ser.} $\leq$)\unsure{$\big/\!16\!\big/$|$\big/\!6\!\big/$} -- lequel \unsure{sera|ser.}

\scanpage{17}\mspage{16}
une relation d'ordre ; en d'autres termes

(i) $\forall x\in\unsure{M|\mathcal{M}}$, dans l'ens ${\unsure{\mathcal{F}|\mathfrak{F}|\mathscr{F}}}^{\unsure{x|X}}$ des $F\in\mathcal{F}$ telles que $\unsure{X|x}\in F$, il y a un plus petit élément, soit $F_{x}$ [NB ce sera l'ens des $Y\in M$ tels que $Y\leq x$)

(ii) L'application $X\mathrel{\text{---}}F_{X}$ est injective, i.e. si $X,Y\in M$, $X\ne Y$, alors $\exists F\in\mathcal{F}$ telle que $X\in F,\ Y\notin F$, ou l'inverse $Y\in F,\ X\notin F$.

Dès lors, les $F\in\mathcal{F}$ \deleted{\unsure{sont des}} \inserted{forment un ens} de parties fermées de $M$, \inserted{contenant les $M_{\leq x}$, $x\in M$\unsure{!?|!|?}} \deleted{\unsure{Se donner un} $\mathcal{F}\subset\unsure{\mathcal{P}|\mathrm{P}|\mathfrak{P}}(M)$ \unsure{satisfais.|satisfais|satisfasse} (i) (ii), \unsure{revient au} $\unsure{=|\overline{\mathrm{m}}}$ \unsure{que de se donner}} \inserted{De plus on suppose} $\varnothing$ (partie vide de $M$) $\in\mathcal{F}$ (ce qui implique que $\mathcal{F}$ n'est \uline{jamais} admissible !).

Dans GF $\overline{\underline{\mathrm{IV}}}$ (Analysis Situs, 1\textsuperscript{ère} mouture) j'ai explicité en long et en large la relation entre la donnée de $(\mathcal{F},\leq)$, satisfais. à des conditions qui y étaient appelées $\underline{\mathrm{C1}}$, $\underline{\mathrm{C2}}$ $\underline{\mathrm{C3}}$ (\unsure{pages|page} \deleted{\unsure{52|53}}53 \deleted{\unsure{56|54}}54, 58),
\margin{J'appellerai ces conditions plutôt $\underline{\mathrm{At}}\,1$, $\underline{\mathrm{At}}\,2$, $\underline{\mathrm{At}}\,3$}
\deleted{que} et un ens ordonné $M$, muni d'un ens de parties fermées $\mathcal{F}$, contenant les $M_{\leq\deleted{\unsure{x|y}}x}$, la partie vide, \inserted{et} avec tout $F$ ses parties fermées (pour l'ordre de $M$), plus la condition : si $F,G\in\unsure{\mathcal{J}|\mathcal{F}}$, tels que $\forall x\in F,\ \unsure{y|\forall y}\in G$ on ait $\unsure{\overset{\scriptscriptstyle\cdots}{\scriptscriptstyle v}|\overset{\scriptscriptstyle\cdots}{\scriptscriptstyle\vee}}\,M_{\leq x}\cup M_{\leq\unsure{y|Y}}\in\mathcal{F}$, alors $F\cup G\in\mathcal{F}$. (critère de compatibilité).

\scanpage{18}
\mspage{17}

Il n'y a pas à y revenir. Je vais plutôt m'attacher à présent aux interprétations en termes d'ombres (grandes ou petites). Il est entendu que dans un premier temps, je m'attache aux \unsure{aspects|aspect} ``algébriques'' ou ``triviaux'' d'un formalisme des figures, et me contente de dégager \unsure{les|le} \unsure{bons|bon} \unsure{axiomes|axiome}, de telle façon que

1\textsuperscript{o}) Ceux-ci \unsure{\deleted{doivent}|doivent} valables dans tout ``atelier'' de la forme $\mathrm{Fig}(\deleted{\unsure{E|L}}L)$, formé des figures ensemblistes dans un ensemble donné $L$ (sans conditions de finitude d'aucune sorte),

et

2\textsuperscript{o}) qu'ils soient applicables aussi, disons, aux figures \uline{simpliciales} affines finies, dans un espace affine sur $\mathbb{R}$ (ou, plus gén\textsuperscript{t}., sur un corps ordonné quelconque).

Pour expliciter, en termes de l'ens ordonné $(\unsure{M|\mathcal{M}},\leq)$, \unsure{les ens|l'ens} $\unsure{\mathcal{F}|\mathfrak{F}|\mathscr{F}}$ des figures, en tant qu'ens de parties de $M$, \deleted{il \unsure{est}} ou du moins des \uline{figures de type fini}
\[
\begin{aligned}
 F\in\mathcal{F}\ \text{de t.f} \quad&\Longleftrightarrow\quad
 \widetilde F=\deleted{\unsure{\text{fibres}|\text{fibre}}}\,\text{réunion finie}\\
 &\hspace{4.9em}\text{d'un}\,M_{\leq \unsure{x_i|X_i}}\\
 &\hspace{3.0em}\text{(pour un nb \uline{fini}}\\
 &\hspace{3.5em}\text{de strates de }\deleted{\unsure{M|x_i|X}}\text{)}
\end{aligned}
\tag*{(1.34)}
\]

\scanpage{19}
\mspage{18}

la chose essentielle est la connaissance de
\[
\underset{\unsure{M|\mathcal{M}}}{\between}\quad\text{\unsure{relt\textsuperscript{n}|relat\textsuperscript{n}|rel\textsuperscript{n}} binaire sym. réfl. (sur }M\text{)}
\tag*{(1.35)}
\]
\deleted{\unsure{qui indique|qui implique}} satisfaisant
\[
\unsure{X|x}\between\unsure{Y|y},\quad X'\leq X,\,Y'\leq Y\quad\Longrightarrow\quad X'\between Y'
\tag*{(1.36)}
\]

\deleted{\unsure{elle|alors}}

il y a corr. biuniv. \unsure{entre|en} l'ens \unsure{des|de} telles relations sur $M$, et l'ens des $\unsure{\mathcal{F}|\mathfrak{F}|\mathscr{F}}\inserted{\subset\unsure{\mathcal{P}|\mathrm{P}|\mathfrak{P}|\mathscr{P}}(M)}$ satisfaisant les conditions énumérées précédemment, plus celle\deleted{s d'être} que les $F\in\mathcal{F}$ soient \deleted{\unsure{toutes|toute|d'être}} de type fini\unsure{.|e|s} C'est en ce sens que j'ai pu dire que ``moralement'', le couple $(M,\mathcal{F}\subset\mathcal{P}(M))$ équivalait à \deleted{\illegible{1}} la donnée de $\underset{M}{\leqslant}$ seul \unsure{de|des} $\between$ satisfaisant (1.36).

\deleted{Mais il s'agit maintenant}

Pratiquement, dans presque tous les exemples, c'est via $M$\inserted{(\unsure{multistrates|multistrate})}, muni de $\leq$ et d'une relation de disjonction $\between$, qu'on d\unsure{é|e}finir\unsure{a|.} les ateliers\ldots

Il faut à présent en venir à la relation $\ll$. \deleted{\unsure{Comme}} Elle va pour l'instant servir surtout à former des ombres, Je commence par les grandes ombres

\scanpage{20}
\mspage{19}

\[
 \mathrm{Omb}(F)=\{X\in\unsure{M|\mathcal{M}}\mid X\ll F\}\,\deleted{\unsure{\subset|=}}
 \tag*{$(*)$}
\]

et la multiombre
\[
 \mathrm{Multom}(F)=\{\mathrm{Omb}(X)\mid X\in\widetilde F\}.
 \tag*{$(**)$}
\]
\margin{$\widetilde F=F$, si on identifi $F$ à une partie de $M$ via $\unsure{\mathcal{F}|\mathfrak{F}|\mathscr{F}}\unsure{\mapsto|\to}\unsure{\mathcal{P}|\mathrm{P}}(M)$\\ $F\unsure{\mapsto|\vdash}\widetilde F$}

En apparence $(*)$ utilise la connaissance de $\ll$ dans $\mathcal{F}$ tout entier, mais \deleted{\illegible{1}}\inserted{il} \deleted{\unsure{on dira qu'elle|voudra qu'elle}\unsure{,}} \inserted{r\unsure{é|e}sultera des axiomes} qu'elle ne dépend que de $\ll$. En tous cas, il est clair que la multiombre ne dépend que de $\underset{M}{\ll}$ (au lieu de $\underset{\mathcal{F}}{\ll}$).

\margin{\uline{NB} Je \unsure{vais|vois} changer plus bas le point de vue \unsure{pris|pris pour|pour} $\ll$ et les axiomes C4, C5.}

Je veux que cette multiombre soit une \uline{figure ensembliste dans} $M$, \uline{ayant pour support l'ombre}, \inserted{et indexée fid\unsure{è|e}lement par $\widetilde F$\,$\unsure{\phi}$} Cela nous amène à
\margin{$\phi$ avec la relation d'ordre $\leq$ dans $\widetilde F$ correspondant à l'inclusion de $\mathrm{Multomb}(F)$, i.e. $\downarrow$}

\inserted{\uline{C4} Si $X,Y\in M$ tels que $X\between Y$ (i.e. $X,Y$ font partie d'une = figure $F$), alors $X\ll Y\Rightarrow X\leq Y$}

\uline{C\deleted{\unsure{4}}5}: Si $X\in M$, $F\in\mathcal{F}$, \inserted{tels que $X\ll F$, et} si $F^X$ est l'ens
\[
 \text{\deleted{\unsure{des|plus}} des }Y\in M\quad\text{tels que }\deleted{\unsure{X\ll Y|X\leq Y}\,\unsure{:}}\,X\ll Y\leq F
 \tag*{(1.37)}
\]
(l'ens des strates de $F$ \unsure{raffinées|raffinée} par $X$), \deleted{\unsure{on}} cet ens $F^X$ . admet un plus petit élément \inserted{pour $\leq$} (i.e. il existe une strate de $F$ raffinée par $X$, et parmi \deleted{\unsure{toutes}} les strates de $F$ raffinées par $X$, il en est une plus petite)

\margin{\uline{Cor} \inserted{Si $X\in M$, $F\in\mathcal{F}$} Pour que $X\ll F$, il f et S qu'il existe $Y\deleted{\in\unsure{\widetilde F|\mathcal{F}}}\in M$ avec $\unsure{X|x}\ll\underbrace{Y\leq F}_{Y\in\widetilde F}$. Ainsi $X\ll F$ \unsure{dans|de} $M\times\mathcal{F}$ connu via $X\ll Y$ \unsure{dans|de} $M\times M$}

\deleted{\uline{C6} Soit $\unsure{F\in\mathcal{F},\,X\in\widetilde F}$ \unsure{puis|soit} $X\in M$, $\unsure{X\ll F|X\in\mathcal{F}}$.}

\deleted{Alors $\exists Y\in\unsure{M|\mathcal{F}M}$ avec \unsure{*|t.q.} a) $Y\ll X$ i.e. $\unsure{X|Y}\in F^Y$ \quad b) $Y\ll X'\leq X$ et $X'\in M\Longrightarrow X'\unsure{=|\simeq}X$}

\deleted{b) \unsure{que|pour} $X$ est \unsure{la|le} plus \uline{petit} \uline{\unsure{élément}} de \unsure{$F$} (cf C4)}

\margin{\uline{NB} Le fait que dans $\mathrm{Multom}(X)$, les $\mathrm{Omb}(X)^{\unsure{\circ|o}}\ne\varnothing$, est trivial\,; car $X\in\mathrm{Omb}(X)^\circ$}

\scanpage{21}
\mspage{20}

Ainsi on trouve donc une application
\[
 \unsure{\mathcal{F}|\mathfrak{F}|\mathscr{F}}\longrightarrow\mathrm{Fig}(\unsure{M|\mathcal{M}}),\quad F\longmapsto\mathrm{Multomb}(F)
 \tag*{$(*)$}
\]
Elle est injective, car \inserted{$X$ est connu quand $\unsure{\mathrm{Om}|\mathrm{Omb}}(X)$ est connu (c'est le plus grand élément de $\mathrm{Omb}(X)$ pour $\ll$), donc $\widetilde F$ est connu quand on \unsure{connaît|connait} $\mathrm{Multom}(F)$, donc $F$ aussi.}

Il est immédiat
\[
 \deleted{F}G\leqslant\deleted{G}F\quad
 \unsure{\Longleftrightarrow|\Longrightarrow}\quad
 \underbrace{
 \underbrace{\mathrm{Multomb}(\deleted{F}G)}_{
 \substack{\text{les strates sont}\\\text{les }\mathrm{Omb}(\deleted{X}Y)\\\deleted{X}Y\in\widetilde{\deleted{F}G}}}
 \leqslant
 \underbrace{\mathrm{Multomb}(\deleted{G}F)}_{
 \substack{\text{les strates}\\\text{sont les}\\\mathrm{Omb}(X)\\X\in\widetilde F}}}_{
 \substack{\text{or }\widetilde{\unsure{G|C}}\subset\underset{\unsure{.}}{\widetilde F}\text{ est}\\
 \text{une partie}\\\text{fermée de }\widetilde F\\\text{si }G\leq F}}
 \tag*{(1.38)}
\]

\uline{19 Juin}\qquad Finalement, j'ai changé de point de vue pour le sens géométrique de la relation $\deleted{F}G\ll\deleted{G}F$, que je concevais comme ``$\unsure{G|C}$ est une sous figure d'une subdivision de $F$''. Déjà dans \deleted{\unsure{les}} le cas simplicial (i.e la structure combinatoire des \unsure{\deleted{tout}|\deleted{tous}|\inserted{$\mapsto$}} multistrates étant \deleted{\unsure{celle|elle}} une structure de simplexe : l'ens. ordonné des sous-ens d'un ens. fini \ldots), cela n'a pas l'air naturel. Par exemple, si on a une théorie modérée, et deux simplexes modérés $X,Y$ dans un $\mathbb{R}^n$ disons, on dira que $\deleted{X}Y\ll X$ si c'est le cas pour les figures ensemblistes correspondantes -- ce qui n'implique sans doute pas l'existence d'une subd. de $X$ telle que $Y$ en soit une multistrate \unsure{de|et} \unsure{la|Ce|ce}

\scanpage{22}
\mspage{21}

question s'il en existe une semblerait accessoire. Aussi, finalement, j'ai envie de poser

\noindent\begin{tabular}{@{}||@{\hspace{.55em}}p{.93\linewidth}@{}}
\uline{A\deleted{C}t 4}\quad Si $\deleted{G}F,\deleted{F}G\in\unsure{\mathcal{F}|\mathfrak{F}|\mathscr{F}}$, alors $\deleted{F}G\ll\deleted{G}F$ sss pour tout $X\in\widetilde{\deleted{F}G}$, $\unsure{\deleted{\exists}|\deleted{s}|\&}$ $\exists Y\in\widetilde F$ avec $X\ll Y$
\end{tabular}

Ainsi, la relation $\ll$ sur $\mathcal{F}$ se déduit \margin{$\times$} de la relation $\ll$ sur $\unsure{M|\mathcal{M}}$. Donc, si on \unsure{suppose|s'impose|dispose} \unsure{les|ces} figures de type fini, $(\mathcal{F},\leq,\ll)$ (i.e. tout l'atelier) se \unsure{réconstitue|réconstituer} entièrement à partir de $(M,\leq,\ll,\between)$.

\margin{\uline{Cor} $\forall F\in\mathcal{F},\ \mathrm{Omb}(F)=\bigl|\unsure{\mathrm{Multomb}|\mathrm{Multiomb}}\,F\bigr|$}

\noindent\begin{tabular}{@{}||@{\hspace{.55em}}p{.93\linewidth}@{}}
\uline{\deleted{C}At 5}\quad Soient $X\in M$, $F\in\mathcal{F}$, tels que $X\ll F$. Alors \deleted{dans} l'ensemble $F^X=\{Y\in\widetilde F\mid\underbrace{X\ll Y}_{\text{i.e. }X\in\mathrm{Omb}(Y)}\}$ (lequel est $\ne\unsure{\varnothing|\phi}$ par C4) admet un plus petit elément\unsure{$\downarrow$|*|,} (pour $\leq$)
\end{tabular}

En d'autres termes, l'ensemble de parties $\mathrm{Multom}\deleted{\unsure{b}}b(F)=\{\widetilde Y\mid Y\in\widetilde F\ \text{i.e.}\ Y\triangleleft F\}$ est une \uline{figure} dans $M$, \deleted{[\uline{NB} \unsure{étant entendu} que \unsure{les|le} $\widetilde Y^{\unsure{\circ|o}}=\widetilde Y\setminus\bigcup_{\unsure{X<|X\ll}}$} ou du moins, est \uuline{pré}\uline{admissible}. C'est en tous cas impliqué par \uline{C\unsure{5|6}}, et lui est équivalent si on admet que la relation (pour $Y,Y'\in\widetilde F$) $Y\leq Y'$, qui implique $Y\ll Y'$ \unsure{d'où|donc}. $\deleted{\unsure{G|O}}\mathrm{Omb}(Y)\subset\mathrm{Omb}(Y')$ lui est en fait équivalente. Cela équivaut à la condition

\noindent\begin{tabular}{@{}||@{\hspace{.55em}}p{.93\linewidth}@{}}
\uline{\deleted{C}At 6}\quad Si $X,Y\in M$ et $X\between Y$, alors $X\ll Y\Rightarrow X\leq Y$ (donc $X\ll Y\ \Longleftrightarrow\ X\leq Y$).

En d'autres termes,
\end{tabular}

\scanpage{23}
\mspage{22}

\uline{Corollaire} 1 Soit $F\in\unsure{\mathcal{F}|\mathfrak{F}|\mathscr{F}}$. Alors dans $\widetilde F\subset\unsure{M|\mathcal{M}}$, les relations d'ordre $\leq$ et $\ll$ coïncident.

\uline{Corollaire 2} Soient $F$ et $G$ deux figures compatibles $(F\between G)$. Alors $F\ll G$ \unsure{sss|ssi} $F\leq G$.

Il suffit de prouver $\Longrightarrow$, et par C4, on est ramené à \deleted{\unsure{prouver que}}\inserted{le voir} pour $X$ ($\in\widetilde F$) et $Y$ ($\in\widetilde G$) dans $M$, ce qui n'est autre que C6.

Ainsi, pour \unsure{tout|toute} $F\in\mathcal{F}$, l'indexation de $\mathrm{Multomb}(F)$ par $\widetilde F$ est fidele, et respecte les relations d'ordre naturelles \inserted{\unsure{habituelle|habituelles} $\leq$}.

Ceci \unsure{\deleted{a}|\deleted{e}|4|+} dit, pour voir que $\mathrm{Multomb}(F)$ est bien une figure, il reste à voir que
\[
 \text{pour }Y\in\widetilde F,\qquad \widetilde Y^{\unsure{\circ|o}}\overset{\mathrm{def}}{=}
 \widetilde Y\setminus
 \bigcup_{\substack{Y'\in\widetilde F\\Y'\unsure{<|\leq}Y\;(\text{ou encore, }Y'\ll Y,\,Y'\ne Y)}}\widetilde Y'\quad\text{est}
\]
non vide. Mais il n'y a qu'à \deleted{\unsure{peut-être|prouver|remarquer}}\inserted{noter} que $Y\in\widetilde Y^\circ$ !

\margin{NB. pourrait remonter \unsure{avant|au cas} C6, \unsure{car|bien} il n'y avait pas besoin de \unsure{savoir|voir} que $Y'\ll Y$, $Y'\ne Y$ équivaut à $Y'\unsure{<|\underset{\cdot}{<}}Y$ pour $Y,Y'\in\widetilde F$ !}

Par la relation d'ordre $\ll$ sur $M$, on trouve une application
\[
 \begin{aligned}
 M&\lhook\joinrel\longrightarrow\unsure{\mathcal{P}|P|\mathfrak{P}|\mathscr{P}}(M)\\
 X&\longmapsto\mathrm{Omb}(X)=M_{\ll X}=\{Y\in M\mid Y\ll X\}
 \end{aligned}
 \tag*{(\deleted{\unsure{*|1.38}}1.39)}
\]
qui est \uline{injection} ($\ll$ relation d'\uline{ordre}, pas seulement préordre) \inserted{et \unsure{la} relation d'ordre de $M$ \unsure{induit|induite} par celle de $\mathcal{P}$}. Les figures $\mathrm{Multom}(F)$ sont les images, par cette application, des $\widetilde F$ ($F\in\mathcal{F}$). \deleted{Si $F$ et $G\in\mathcal{F}$,}

\scanpage{24}
\mspage{23}

\uline{Proposition} Soient $F,G\in\unsure{\mathcal F|\mathfrak F|\mathscr F}$. Conditions équivalentes
\[
\begin{aligned}
&\text{(i)}\quad F\leq G\\
&\text{(i bis)}\quad \unsure{\widetilde F|\widetilde{\widetilde F}}\subset\widetilde G\\
&\text{(ii)}\quad \mathrm{Mul\unsure{t|\deleted{t}}omb}(F)\leq\mathrm{Mul\unsure{t|\deleted{t}}omb}(G)
 \qquad\begin{gathered}\text{i.e. }\mathrm{Multomb}(F)\\\text{est une sous-figure}\\\text{de }\mathrm{Multomb}(G)\end{gathered}\\
&\text{(ii bis)}\quad \mathrm{Mul\unsure{t|\deleted{t}}omb}(F)\subset\mathrm{Mul\unsure{t|\deleted{t}}omb}\unsure{G)|(G)}
\end{aligned}
\]

Dém. $F\leq G$ \deleted{trivial, et} $\unsure{\Longleftrightarrow|\Longleftarrow}\widetilde F\unsure{\subset|\in}\widetilde G$ \deleted{trivial, et}\inserted{déjà connu} \deleted{\unsure{l'inverse}}\inserted{et mis} pour mémoire ($\Longrightarrow$ trivial, $\Longleftarrow$ car $\deleted{\widetilde F}F=\mathop{\mathrm{Sup}}\limits_{X\in\widetilde F}X$, $G=\mathop{\mathrm{Sup}}\limits_{Y\in\unsure{\widetilde G|\overline G}}Y$). L'équivalence (i bis) $\Longleftrightarrow$ (ii bis) triviale par l'injectivité de (\deleted{\unsure{*|1.38}}1.39). Or (ii bis) est impliqué par (ii) (tautologique) \unsure{prouvons|poussons} l'inverse, donc que (ii bis) implique que $\mathrm{Mulomb}(F)$ est fermé dans $\mathrm{Mulom}\,G$ pour l'inclusion. Mais \unsure{on aura} $\widetilde F\subset\widetilde G$ \deleted{et il \unsure{est}}\inserted{donc} \deleted{et \unsure{on|en}} $F\leq\unsure{G|C}$, et il est \deleted{\unsure{connu}}\inserted{connu et trivial} que $\widetilde F$ est \inserted{alors} fermé dans $\widetilde G$.

\uline{Corollaire 1} \deleted{Conditions équivalentes\unsure{,} pour}\inserted{Soient} $F,G\in\unsure{M|\mathcal M}$ \inserted{tels que $F\unsure{\between}G$, alors} \deleted{(i) $F\between G$}
\[
 \deleted{\text{(ii)}}\quad \mathrm{Mulomb}(F)\between\mathrm{Mulomb}(G),
 \tag*{(*)}
\]
\deleted{Quand elles sont \unsure{remplies|vérifiées|compatibles},}\inserted{et} on a
\[
 \left\{\begin{aligned}
  \mathrm{Mulomb}(F\vee G)&=\mathrm{Mulom}(F)\cup\mathrm{Mulomb}(G)\\
  \mathrm{Mulomb}(F\wedge G)&=\mathrm{Mulom\unsure{b|s}}(F)\cap\mathrm{Mulomb}(G).
 \end{aligned}\right.
 \tag*{(**)}
\]

\deleted{La condition (ii) signifie que $\widetilde F$ et $\widetilde G$ sont fermés dans $\widetilde F\cup\widetilde G\subset M$, pour la relation $\ll$ (\unsure{induite|induit} par $\subset$ dans $\unsure{\mathcal P|\mathfrak P|\mathscr P|P}(M)$).}

\scanpage{25}
\mspage{24}

En effet, $\widetilde{F\vee G}=\widetilde F\cup\widetilde G$, et $\widetilde F$ et $\widetilde G$ y sont fermés, \deleted{\unsure{dans|de} $\unsure{M|\mathcal{M}}$} ce qui prouve que
\[
 \mathrm{Multomb}(F\cup G)=\mathrm{Multomb}(F)\cup\mathrm{Mulomb}(\unsure{C|G}),
\]
et que $\mathrm{Mulomb}(F)$ et $\mathrm{Mulomb}(G)$ y sont fermés, donc \unsure{(*)} et la première relation \unsure{(**)|($\ll$)}. La deuxième relation \inserted{\unsure{(**)|($\ll$)}} provient de $\unsure{=|\eqsim|\overline{m}}$ de $\widetilde{F\cap G}=\widetilde F\cap\widetilde G$.

\uline{NB} Il n'est pas vrai en général que $\mathrm{Mulomb}(F)\between\mathrm{Mulomb}(G)$ implique $F\between\unsure{C|G}$, même si $F$ et $G$ sont \unsure{réduits|s'écrivent} aux lieux distincts $x$ et $y$ --- auquel cas $\mathrm{Multomb}(X)=\{\{\unsure{\lambda|x}\}\}$, $\mathrm{Mulomb}(Y)=\{\{y\}\}$ sont compatibles et $\unsure{=|\eqsim}$ \unsure{disjointes|disjoints}, mais $x$ et $y$ ne le sont pas forcément. Mais il \unsure{ser.|sera} vrai \unsure{par contre qu'on a équivalence|pour d'autres quasi-équivalences|pour d'autre quasi-équivalence}, dans le cas où $\unsure{\mathcal{F}|\mathfrak{F}|\mathscr{F}}$ est \unsure{un|une} atelier \unsure{ensembliste|ens.embliste}\unsure{.|:}
\margin{\unsure{vérifier|vérifies} plus tard}

\uline{Corollaire 2} Si $F$ et $G$ sont \unsure{disjoints|disjointes|disjoint}, alors $\unsure{\mathrm{Mul.mb}|\mathrm{Mulomb}}(F)$ et $\mathrm{Mulomb}(G)$ aussi.

Cela provient formellement du corollaire 1 et de la formule
\[
 \mathrm{Mulomb}(\unsure{\varnothing|\phi}_{\unsure{\mathcal{F}|F}})=\varnothing\quad(\text{figure vide dans }\unsure{M|\mathcal M})
 \tag*{(1.40)}
\]

\uline{Proposition} Soient $\unsure{F|\mathbb{F}},G\in\mathcal{F}$. Alors $F\ll G$
\[
 \Longleftrightarrow\quad\mathrm{Mulomb}(F)\ll\mathrm{Mulomb}(G).
\]

\scanpage{26}
\mspage{25}
D'abord une observation. Soit $F\in\unsure{\mathcal{F}|\mathfrak{F}|\mathscr{F}}$, d'où
\[
\mathrm{Omb}(F)=\lvert\mathrm{Multomb}\,F\rvert
 \overset{\mathrm{def}}{=}
 \bigcup_{X\in\widetilde F}\deleted{\unsure{\widetilde X|X|x}}\mathrm{Omb}(\unsure{X|x})
\]
On a une application canonique
\[
\mathrm{Omb}(F)\xrightarrow{\varphi_F}\widetilde F
\]
\margin{NB. $\varphi_F$ induit $\mathrm{id}_{\widetilde F}$ sur $\widetilde F\subset\mathrm{Omb}(F)$}
associant à tt $Y\in\mathrm{Omb}\,F$, i.e. $Y\in\unsure{M|\mathcal{M}}$ tel que
$\deleted{\unsure{Y\ll F|Y\in F}},\quad Y\ll F$, le plus petit élément de $\deleted{\unsure{G}}F^{Y}$
i.e. le plus petit $X$ de $\deleted{\unsure{G}}\widetilde F$ parmi ceux \unsure{satisfaisant}
$Y\ll X$ i.e. $\deleted{\unsure{O|\mathrm{Om}}}\,Y\in\mathrm{Omb}\,X$. Cette application est celle qui définit la relation de préordre de $\lvert\mathrm{Mulomb}(F)\rvert$ associée à la figure $\mathrm{Mulomb}(F)$.

Ceci dit, \inserted{pour $\deleted{\unsure{G}}F,\unsure{G|C}\in\mathcal{F}$ donnés,} dire que
$\underbrace{\deleted{\unsure{M}}\mathrm{Mulomb}\bigl(\deleted{F}G\bigr)}_{\Psi}\ll\underbrace{\mathrm{Mulom}\,F}_{\Phi}$,
\unsure{signifie} que \deleted{\unsure{il|elle}}\unsure{l'on|on} a une inclusion
\[
 \lvert\Psi\rvert\subset\lvert\Phi\rvert\qquad\text{i.e.}\quad \mathrm{Om}\,G\subset\unsure{\mathrm{Om}|\mathrm{Omb}}\,F,
\]
et une application \uline{croissante}
\[
\deleted{\unsure{\text{(1.42)}|\text{(1.47)}}}\qquad
f:\deleted{\widetilde F}\widetilde G\mathbin{-}\deleted{\unsure{\widetilde G}}\widetilde F
\]
rendant commutatif le diagramme
\diagram{\[
\begin{tikzcd}[ampersand replacement=\&,column sep=large]
 \mathrm{Omb}\,\deleted{\unsure{F}}G \arrow[r,hook] \arrow[d,"\varphi_{\unsure{G|C}}"] \& \mathrm{Om}\,\deleted{\unsure{G}}F \arrow[d,"\varphi_F"]\\
 \deleted{\unsure{\widetilde F}}\widetilde G \arrow[r,"f" description] \& \widetilde F
\end{tikzcd}\quad .
\tag*{(1.48)}
\]}

\scanpage{27}
\mspage{26}
Montrons que cela équivaut à $G\ll F$. En effet, si $G\ll F$, associant à tout $X\in\widetilde G$ le plus petit élément $\deleted{\unsure{\varphi}}f(\unsure{x|X})$ de $\widetilde F$ tel que
\[
 X\ll f(\unsure{x|X}),
\]
on trouve bien une application
\[
 \deleted{\unsure{\widetilde G|G}}\qquad f:\widetilde G\longrightarrow\widetilde F
\]
et elle est \unsure{évidemment} croissante\,: \margin{si $X,X'\in\widetilde G$}
\[
\begin{aligned}
 X'\leq X\quad\text{i-e.}\quad X'\ll X\inserted{(\ll f(x))}\ &\text{implique}\quad X'\ll f(x)\\
 &\text{donc}\quad\deleted{\unsure{f|\varphi}}\,f(x')\ll f(x)\\
 &\text{qui signifie}\quad f(x')\leq\unsure{f(x)|\deleted{f(x)}}
\end{aligned}
\]
et il faut vérifier la commutativité de (1,42), ce qui \deleted{\unsure{se voit comme} \unsure{;|à} l'instant.} \inserted{signifie}

\uline{Corollaire} Soient $X\in\unsure{M|\mathcal M}$, $F,\unsure{C|G}\in\unsure{\mathcal F|\mathfrak F|\mathscr F}$ tels que $X\ll G\ll F$.
Soit $\deleted{\unsure{X|G|X_G}}\inserted{X'}$ le plus petit élément de $G^X$, $X''$ le plus petit élément de $F^{X'}$, alors $X''$ est \unsure{aussitt|aussi tt|aussi} le plus petit élément de $F^X$.

Il semble qu'il faille le poser en axiome, sous forme en apparence \unsure{affaiblie}\unsure{;|,} mais en fait équivalente. On va \deleted{\unsure{écrire|é}} \inserted{introduire} encore une notation, pour $X,Y\in M$
\[
\begin{aligned}
 X\overset{\circ}{\ll}Y\qquad &(\unsure{-X|\deleted{X}}\text{ est \uline{\unsure{suffisamment} intérieur}}\\
 &\hspace{6em}\text{de }Y)
\end{aligned}
\tag*{(1.48)}
\]
pour
\[
 \deleted{X\ll Y,\ \text{et}\ X\ll Y}
\]
\[
 X\in\mathrm{Omb}(Y)^{\circ}=\mathrm{Omb}(Y)\setminus\bigcup_{Y'<Y}\mathrm{Omb}(Y')
\]
i.e. $X\ll Y$, et $X\ll Y'\leq Y$ implique $Y'=Y$.

\scanpage{28}
\mspage{27}
On doit donc poser

\uline{A\deleted{\unsure{C}}\unsure{t|E|C}7}\quad Si $X,Y,Z\in\unsure{M|\mathcal{M}}$ sont tels que
\[
 X\overset{\circ}{\ll}Y\overset{\circ}{\ll}Z\qquad\text{alors}\qquad X\overset{\circ}{\ll}Z
\]
i.e. la relation $\overset{\circ}{\ll}$ est \uline{transitive}.

Ceci \unsure{assure|conserve} donc $\Rightarrow$ dans la prop.
En sens inverse, S'il \unsure{existe|exist} $f$ comme dans (1.42), soit $X\in\widetilde G$, \deleted{alors $\varphi_G(\unsure{X|x|/}$} \inserted{\deleted{\unsure{on|ce}}} de sorte que $\varphi_G(\unsure{X|x})=X$ (tautologique\deleted{\unsure{?|!}}), moyennant (6), \deleted{on a alors} soit $f(X)=Y$, on a alors $Y=\varphi_F(X)$ i.e. $Y$ est le plus petit élément de $\widetilde F$ \deleted{\unsure{contenant} $X$}. \deleted{\unsure{Donc}} tel que $X\ll Y$. Cela \deleted{signifie} implique donc que pour tt $X\in\deleted{\widetilde F}\widetilde G$, \unsure{existe|exist} $Y\in\deleted{\unsure{\widetilde G|\widetilde C}}\widetilde F$ telle que $X\ll Y$, i.e. $\deleted{\unsure{\widetilde G|\widetilde F}}$ (C4) $\widetilde G\ll\widetilde F$\unsure{..|:}, qed. On n'a pas \unsure{utilisé|utilisé-|utilis-} en fait que $f$ soit croissante, donc

\uline{Corollaire \unsure{à|2} la prop.}\quad \deleted{La condition,} \inserted{Pour} que $F\ll G$ i.e. $\mathrm{Mulom}(F)\ll\mathrm{Mulomb}(G)$. il s. que $\mathrm{Omb}\,F\subset\mathrm{Omb}\,G$,
\margin{Ce ne sera sans doute plus vrai, en \unsure{passant} aux \uline{petites} multifigures}
et que \deleted{\unsure{toute classe}} la relation d'équivalence \inserted{$\unsure{\mathcal R_F|R_F|Q_F}$} sur $\mathrm{Omb}\,F$ définie par $\varphi_F$ soit plus fine que la relation d'équiv. induite \unsure{par|sur} $\mathrm{Omb}\,G$.

\scanpage{29}
\mspage{28}
\uline{Scholie}\quad Ainsi, via $F\longmapsto\mathrm{Mulomb}(F)$,
$(\unsure{\mathcal F|\mathfrak F|\mathscr F},\deleted{\unsure{\ll}}\leq,\ll)$ s'interprète comme un ``sous-atelier'' de $\mathrm{Fig}(\unsure{M|\mathcal M})$.

Donc la donnée de l'atelier $(\mathcal F,\leq,\ll)$ équivaut à celle d'un ens $M$, et d'une partie
\[
 \deleted{\unsure{\Phi|\mathcal F}}\mathcal F_M\,
 \deleted{\unsure{\subset\mathcal P|\mathcal F}}\,
 \deleted{\Phi}\subset\mathrm{Fig}(M)\subset
 \unsure{\mathcal P|\mathrm P|\mathfrak P|\mathscr P}(\mathcal P(M))
 \tag*{(1.44)}
\]
\[
 \left(\text{\unsure{donc|d.---\illegible{1}}}\quad
 \mathcal F_M\,\deleted{\unsure{\Phi}\,\unsure{\mathcal F_M|\Phi_M}}
 \in\mathcal P(\mathcal P(\mathcal P(M)))\right)
\]
satisfaisant certaines conditions \inserted{(qui doivent traduire C1 ; C7)}.

Du coup, la notion de \inserted{(grande)} multiombre reprend un plus grand intérêt, et je me demande s'il ne faudrait pas une notation \deleted{\unsure{plus}} et un nom \unsure{plus|pluss} lapidaires.

Je voudrais expliciter les conditions sur cet ens de figures ensemblistes dans $M$\unsure{.|-} Je désignerais par $\mathrm{At}^{*}1$ etc les axiomes en termes de $(M,\deleted{\unsure{\Phi_M}}\mathcal F_M\in\mathcal P\mathcal P\mathcal P(\unsure{M|m}))$.
\deleted{\unsure{$(\Phi_1)$|$(\Phi_F)$|$(\Phi_{\mathcal F})$|(1.49)|(1.50)}} Il est entendu que $\deleted{\unsure{\Phi}}\mathcal F_M\subset\mathrm{Fig}(M)$, i.e. \deleted{\uline{\unsure{At|AE}$^{*}$1}} que les éléments de $\deleted{\unsure{\Phi}}\mathcal F_M$ sont des ensembles de parties de $M$ qui sont des \uline{figures} ensemblistes $F$ dans $M$. \deleted{\illegible{1}} \deleted{Ainsi} Considérons

\scanpage{30}
\mspage{29}
\[
\begin{aligned}
 \overset{\deleted{\unsure{\sim|\approx}}}{\unsure{M|\mathcal M}}^{*}
  &=\bigcup_{\deleted{\unsure{F}}\Phi\in\deleted{\unsure{\Phi}}\unsure{\mathcal F|\mathfrak F|\mathscr F}_M}
       \deleted{\unsure{F}}\Phi\subset\unsure{\mathcal P|P|\mathfrak P|\mathscr P}(\deleted{\unsure{F}}M)\\
 &\quad\unsure{\shortparallel|\text{''}}\\
 &\left\{\unsure{A|\dot A}\in\mathcal P(M)\mid\exists\,\deleted{\unsure{F}}\Phi\in\deleted{\unsure{\Phi}}\mathcal F_M
       \ \text{\unsure{avec}}\ A\in\deleted{\unsure{F}}\Phi\right\}
\end{aligned}
\]
i.e. l'ensemble des strates de toutes les figures $\in\deleted{\unsure{\Phi}}.\mathcal F_M$

\uline{At$^{*}$1}$_{\unsure{0|o}}$\quad $\overset{\deleted{\sim}}{M}^{*}$ est une figure ensembliste
\margin{\rule{0.45pt}{3.5\baselineskip}}
dans $M$, de support $M$. De plus, le préordre associé de $M$ \inserted{(\unsure{soit|sur} $\ll$)} est une relation d'\uline{ordre} i.e.
\[
\begin{aligned}
 \unsure{x|X},\unsure{y|Y}\in M,\ \unsure{x|X}\neq\unsure{y|Y}\quad\Longrightarrow\quad \exists\, A\in\overset{\deleted{\sim}}{M}^{*}
 &\quad\text{\unsure{avec}}\quad X\in A,\ Y\notin A\\
 &\quad\text{ou}\qquad Y\in A,\ x\notin A
\end{aligned}
\]

Ainsi, l'application
\[
 \unsure{x|X}\longmapsto \deleted{\unsure{M^{*}_{\ll x}|\widetilde M_{\ll x}}}M_{\ll\unsure{x|X}}
 \overset{\mathrm{def}}{=}\mathrm{Omb}(X)
\]
\[
 \underset{\unsure{\cdot\,\cdot}}{M}\longrightarrow\mathcal P(M)\tag*{(*)}
\]
induit une \uline{bijection} de $M$ \unsure{sur|avec} $\overset{\deleted{\sim}}{M}^{*}$, transformant $\ll$ en la relation d'inclusion.
Les $A\in\overset{\deleted{\sim}}{M}^{*}$ sont les parties fermées \uline{irréductibles} de $M$ admettant un pt générique.

\deleted{\uline{\unsure{At|Al}$^{*}$}2}\quad Soit $\overset{\deleted{\sim}}{\deleted{\unsure{F}}\Phi}\in\deleted{\unsure{\Phi}}\inserted{\mathcal F_M}$, [\unsure{donc|\deleted{donc}} $\overset{\deleted{\sim}}{F}\subset\mathcal P(M)$ et par définition même $\deleted{\unsure{F}}\Phi\subset\overset{\deleted{\sim}}{M}^{*}\subset\mathcal P(M)$), donc $\Phi$ est

\scanpage{31}
\mspage{30}
l'image par \deleted{\unsure{(1.79)|(1.39)|(1.49)}} ($\unsure{*|\alpha}$)
\[
 \unsure{M|\mathcal M}\xrightarrow{\sim}\overset{\deleted{\unsure{\sim|\approx}}}{M}^{*}\subset\unsure{\mathcal P|\mathscr P|\mathfrak P|P}(M)
\]
d'une unique \inserted{partie} $F\subset M$, de sorte que $\Phi=\widetilde F$, i.e.
\[
 \deleted{\unsure{\mathrm{At}|\mathrm{AC}}}\qquad
 \Phi=\bigl\{\underset{\inserted{\substack{\shortparallel\\ M_{\ll x}}}}
 {\mathrm{Omb}(\deleted{\unsure{\widetilde X}}X)}\mid X\in F\bigr\}
\]
Ainsi, on peut aussi interpréter, grâce à $\ll$, $\unsure{\mathcal F|\mathfrak F|\mathscr F}_M$ en termes d'un\deleted{e} \deleted{partie} ens de parties \inserted{$\widetilde{\mathcal F}$} de $M$
\diagram{\begin{tikzcd}[ampersand replacement=\&,column sep=large,row sep=1pt]
 \mathcal F_M \arrow[r,"\sim"] \& \widetilde{\mathcal F}\\
 \cap \& \cap\\
 \mathcal P(\mathcal P(M) \& \mathcal P(M)
\end{tikzcd}}

% slanted
\margin{cf GF $\overline{\underline{\mathrm{IV}}}$ page 54 (Cor. 2) et page 55 (Remarque) Pour C3, cf page 58}
Ceci dit, \deleted{on demande que} \inserted{\unsure{Je veux|on veut|on veux} que $\widetilde{\mathcal F}\subset\mathcal P(M)$,}
\deleted{\unsure{soit associé à un}} \deleted{\unsure{satisfasse}} \inserted{\unsure{pour|muni|muni de}} \deleted{\unsure{cela|une|cette} relation d} la relation d'inclusion, satisfasse C1, C2, C3,
et \deleted{\unsure{identifiant|identifie} les \unsure{à}} pour $M$ \unsure{s'identifiant|s'identifie|s'identifier} à l'ens de ses objets \inserted{irréd.}

\deleted{\uline{\unsure{At|AE}}$^{*}$2 \unsure{bis}}\quad \inserted{On voudra d'abord que} L'ens. \unsure{des|de} parties $\widetilde{\mathcal F}$ \deleted{de} $\unsure{\subset}\mathcal P(M)$ soit \uline{préadmissible}, i.e. pour tout $X\in M$, il existe un plus petit élément de $\widetilde{\mathcal F}$ qui le contienne. En d'autres termes,

\uline{At$^{*}$2}$_{\unsure{0|o}}$\quad
\margin{\rule{0.45pt}{4\baselineskip}}
\inserted{Pour} tout $\deleted{\Phi\in\mathcal F_M}$ $M$-figure ensembliste $\Phi$ \deleted{\unsure{de|dans}} $\in\mathcal F_M$, et pour $A$ [$=\mathrm{Omb}(x)$!] \deleted{\unsure{de}} strate de $\Phi$ (i.e. $A\in\Phi$), il existe une \uline{plus petite figure} \inserted{$\Phi_A$} \deleted{\unsure{tous|dans} \unsure{des}} $\in\overset{\deleted{\sim}}{\mathcal F}_M$
% slanted
\margin{c'est la \uline{\unsure{multistrate}} définie par $A$}
(plus petite au sens de l'\uline{inclusion} \deleted{\unsure{)}} entre figures i.e. entre parties de $\mathcal P(M)$, ou \deleted{de} $M^{*}\subset\mathcal P(M)$) qui contienne $A$ comme élément (laquelle $\Phi_A$ sera donc $\subset\Phi$).

% slanted
\margin{C'est C1 pour $\unsure{\mathcal F_M|F_M}$}
\uline{At$^{*}$}3$_0$\quad Si $\Phi\in\mathcal F_M$, et si \deleted{\unsure{des}} $\Phi_i\in\mathcal F_M$, $\Phi_i\subset\Phi$, alors $\bigcup_{\unsure{i}}\Phi_i\in\mathcal F_M$ (stabilité par réunions \uline{conjointes} \inserted{dans $\mathcal F_M$})

\scanpage{32}
\mspage{31}
et de plus on a $\unsure{\varnothing|\phi}\in\unsure{\mathcal F|\mathfrak F|\mathscr F}_{\unsure{M|\mathcal{M}}}$ (ou encore $\mathcal F_M\neq\varnothing$).

Enfin, on voudra aussi

\uline{At$^{*}$4}$_{\unsure{0|o}}$\quad Si $\Phi,\Psi\in\deleted{\unsure{\Phi}}\mathcal F_M$, pour que $\Phi\cup\Psi$ soit \deleted{élément} dans $\mathcal F_M$ \inserted{(i.e., par At$^{*}$3, $\Phi,\Psi$ \unsure{majorés} dans $\mathcal F_M$ \inserted{pour l'inclusion})}, il f. et s. que pour tout $A\in\Phi$, $\unsure{\cdot|.}B\in\Psi$, $\Phi_A\cup\Phi_B$ soit dans $\mathcal F_M$.

Je voudrais revenir sur ce $\Phi_A$. En fait, $\Phi_A$ consiste en l'ensemble des sous-strates de $A$ dans $\Phi$, donc \uline{At$^{*}$2}$_0$ peut se formuler sous la forme plus compréhensible

\uline{At$^{*}$}2$_0$\quad Soient $\Phi$ et $\Psi$ d\deleted{\unsure{ans}}\unsure{eux} $M$-figures $\in\mathcal F_M$, et soit $A\in\Phi\cap\Psi$ une strate commune. Alors la \unsure{multistr.te|multistrate} associée \unsure{à|:} $A$ dans $\Phi$, et dans $\Psi$, est la \deleted{\unsure{m}}\unsure{même}, et elle est $\in\unsure{,|.}\mathcal F_M$.

\uline{Corollaire}\quad Si $\Phi,\Psi\in\mathcal F_M$ avec $\Phi\subset\Psi$, alors $\Phi$ est \unsure{fermée|fermé} dans $\Psi$, donc est une sous figure $(\Phi\leq\Psi)$.

\scanpage{33}
\mspage{32}
[Il s'ensuit qu'on a \unsure{dans|donc} $\unsure{M|\mathcal M}\simeq M^{*}$ une relation d'ordre $\leq$, avec
\[
\begin{aligned}
 X\leq Y\quad\Longleftrightarrow\quad&
 \deleted{\mathrm{Omb}(x)}\Phi_{\mathrm{Omb}(x)}
 \subset\Phi_{\mathrm{Omb}(Y)}\\
 &\qquad\deleted{\text{i.e.}}\ \text{i.e.}\unsure{,|.}\quad
 \mathrm{Omb}(x)\in
 \underset{\inserted{\substack{\mathrm{def}\ \unsure{\shortparallel|\text{''}}\\
 \unsure{\mathrm{Mulomb}|\mathrm{Multomb}}(Y)}}}{\Phi_{\mathrm{Omb}(Y)}},
\end{aligned}
\]
de telle façon que les \deleted{éléments de} $F\in\widetilde{\unsure{\mathcal F|\mathfrak F|\mathscr F}}$ soient des parties fermées de $M$\unsure{,|.}
Associant à tout $X\in M$ la partie $\deleted{\unsure{\widetilde X|X}}$
\[
\begin{aligned}
 \widetilde X&=\{Y\in M\mid Y\leq\unsure{x|X}\ \text{i.e.}\quad
 \mathrm{Omb}(Y)\in\Phi_{\mathrm{Omb}(x)}\}\\
 &=M_{\leq x}
\end{aligned}
\]
on trouve une application injective
\[
 \overset{\unsure{\cdot\,\cdot}}{M}\hookrightarrow\widetilde{\mathcal F}\quad(\simeq\mathcal F_M)
\]
qui induit un iso entre $M$ et l'ens des éléments str. irréductibles de $\widetilde{\mathcal F}$ (ordonné par inclusion, notée $\leq$).]

\deleted{\unsure{Donc|D}}

\uline{Théorème} - \unsure{\uline{Scholie}|\uuline{Scholie}}\quad La catégorie \inserted{(groupoïde)} des \deleted{\unsure{$\mathcal F$}} ateliers $(\mathcal F,\leq,\ll)$ satisfaisant à \deleted{\unsure{C}}At 1 \unsure{.|à|,} At 7, équivaut à celle des $(M,\mathcal F_M)$ (avec $\mathcal F_M\in\deleted{\unsure{\mathcal P|\mathfrak P|\mathscr P|P}\mathcal P}\mathrm{Fig}(M)$) satisfaisant les conditions suivantes

\uline{At$^{*}$1}, $\mathcal F_M$ \deleted{non vide, et} \inserted{contient la figure vide, et} stable par \deleted{\unsure{les sup majorés|sup majorés}} \inserted{réunions} majorées (dans $\mathcal F_M$)

\deleted{\unsure{Cela|Ceci}} [La stabilité par sup \unsure{équivaudra|équivaudre}, moyennant la condition suivante, à celle de \deleted{\unsure{passage}} stabilité par passage à une \uline{sous} \deleted{\unsure{famille fermée|famille finie}} \uline{figure}.)
\endgroup
